\section{Introduction}
\label{sec:intro}

Estimating the 6D pose of a known object from a single RGB image is ambiguous when the visible part of the object looks the same from
several viewpoints. Symmetry causes this, as do occluded distinguishing features and texture-less surfaces that offer little to
disambiguate~\cite{hodan2017tless,hodan2020epos,murphy2021ipdf}. Treating the pose as a distribution rather than a single value handles
the ambiguity, and diffusion models on SE(3) have become an effective way to represent such distributions.
LiePose~\cite{hsiao2024liepose} samples poses with a score-based diffusion process on SE(3) and reaches high accuracy on T-LESS
without symmetry annotations.

A diffusion denoiser, however, is an open-loop function. At each step it maps image features and the current noisy pose to a
correction, and it never checks whether the current hypothesis is consistent with the image. That check is the core of
render-and-compare refinement~\cite{li2018deepim,labbe2020cosypose,labbe2022megapose,wen2024foundationpose}: render the object at the
hypothesis, compare the rendering with the observation, and update. Render-and-compare has been the most accurate final stage in BOP
for years~\cite{hodan2020bop,sundermeyer2023bop22}. It requires a renderer and a textured 3D model, however, runs as a separate refiner
after an initial estimate, and does not by itself represent a multi-modal distribution.

We ask whether the comparison can be moved into the diffusion process, in a learned feature space, so that every denoising step of
every sample is conditioned on how well its hypothesis explains the image. This needs a model that answers ``what would the image
features be if the object were at pose $g$?'' for arbitrary $g$, cheaply enough to be queried at every step. We build such a model as a
world model in the sense of predictive latent models~\cite{assran2023ijepa,bardes2024vjepa,zhou2025dinowm}, with the object pose as
its action. It operates on frozen DINOv2~\cite{oquab2024dinov2} patch tokens, and it is a \emph{feature field}: tri-plane
features~\cite{chan2022eg3d} and a small MLP give, at every point in a box around the object, a density and a latent vector. The field
renders the tokens at a pose by volume rendering~\cite{mildenhall2020nerf} rays through the observation's crop, and it is fitted
once per dataset to the tokens of posed training crops. Everything it renders is a DINOv2 token or a latent. It has no RGB output and is
given no surface, mesh or texture; its geometry is whatever density best explains the tokens. Because the field renders a pose rather
than borrowing a nearby training view, a small rotation of the hypothesis moves every ray by exactly that rotation. Its prediction
carries no lighting, background or occlusion of another image.

Our method, \method (Latent Render-and-Compare), is a diffusion-based pose estimator built around this world model
(\cref{fig:teaser}). It replaces the usual image-encoder-plus-pose-MLP denoiser with a closed loop. At every step the field renders,
for each sample's current pose, the tokens in the observation's crop and the field's latent at the surface each ray hits. A token
transformer reads the observed tokens, the rendered tokens, their difference and the latent, conditioned on the noise level and the
pose. The image encoder is a frozen foundation model, and the object's appearance as a function of pose is held by the frozen field
rather than learned into the denoiser. The same rendering defines an energy, the foreground-weighted distance to the observation, which
ranks a set of samples at test time without ground truth. We compare against a control that is identical to \method except that it
receives no rendered quantity, so the contribution of the comparison is measured in isolation.

We evaluate on T-LESS~\cite{hodan2017tless} under the BOP protocol~\cite{hodan2018bop} and on the SPEED+ spacecraft
benchmark~\cite{park2022speedplus} with the official SPEC2021 score, against LiePose and the control trained by us with the same data
and crops. We also report what the field can discriminate on its own, and compare it with a retrieval-referenced world model that
predicts tokens from the nearest training view (\cref{sec:supp-ret}). \TODO{summarise the main findings in one paragraph once
Tabs.~\ref{tab:bop}--\ref{tab:speedplus} are filled; no claim before the numbers exist.}

\noindent Our contributions are:
\begin{itemize}
\item A feature-field world model for object pose. It renders frozen DINOv2 tokens at an arbitrary pose by volume rendering a
tri-plane field, is fitted once from posed training crops, and needs no reference view or textured model (\cref{sec:wm}).
\item \method, a closed-loop denoising architecture for SE(3) diffusion. A frozen foundation-model encoder and a token transformer
compare the field's rendering at the current hypothesis with the observation at every step. The same rendering gives an energy for
selecting among samples without ground truth (\cref{sec:den,sec:select}).
\item A controlled evaluation against LiePose at a matched training budget and a no-world-model control, with ablations of what the
denoiser reads from the field and an analysis of what the field can and cannot discriminate (\cref{sec:exp}).
\end{itemize}
